\section{Square-root neighbours, Cheon, and cross-validation}
\label{sec:rho-variants}

Rho is the reference. The neighbours are priced in the same unit so
that a later claim has somewhere to sit.

\subsection{Two grumpy giants and a baby}

Bernstein--Lange~\cite{bernsteinlange2012grumpy} interleave three walks
after shifting so that $x'=x-\ell$ lies in $[0,N)$:
\[
  B_j = jG,\qquad
  U_i = H' + imG,\qquad
  V_i = 2H' - i(m+1)G.
\]
Collisions $B_j=U_i$, $B_j=V_i$, $U_i=V_k$ each give a congruence for
$x'$; all are taken modulo the group order when it is known, and every
candidate is verified. The step $m=\alpha\sqrt{N}$ was chosen by
simulating the three walks in exponent space and confirming with the
implementation (400 instances): $\alpha=0.7$ gives $S=1.18$ on the
whole group and $S=1.78$ on a $2^{22}$ interval without wrap. The
paper's success probability after $1.5\sqrt{n}$ operations is
$0.71875$; the simulation gives $0.70$ at $n=4093$. Interleaved BSGS
in the same simulation is $1.33$, the known $4/3$
figure.\art{docs/ALGORITHMS.md, grumpy table}

\subsection{Free precomputation}

Bernstein--Lange free precomputation~\cite{bernsteinlange2012small,bernsteinlange2013}
builds, once per group and base, a table of distinguished-point chains
and then solves each subsequent target in far fewer operations than a
from-scratch square-root search. With $t=\mathrm{round}(\log_2 n/3)$,
the table has $\sim n^{1/3}$ chains, the build costs $\sim n^{2/3}$
group operations, and each online target costs $\sim n^{1/3}$.
Table~\ref{tab:precomp} is the one-core measurement.\art{docs/BENCHMARKS.md, ``Discrete logarithm with precomputation''}

\begin{table}[t]
\centering
\small
\caption{Bernstein--Lange precomputation, one core. $P$ is the
  one-time build; $T$ the per-target online cost. $\sqrt{n}/T$ is the
  online speed-up over a from-scratch search.}
\label{tab:precomp}
\begin{tabular}{@{}lrrrrrr@{}}
\toprule
grp & bits & $t$ & chains & $P/n^{2/3}$ & $T/n^{1/3}$ & $\sqrt{n}/T$ \\
\midrule
$\ZZ_p^*$ & 24 & 7 & 200 & 1.398 & 1.926 & 6.6 \\
$E(\FF_p)$ & 24 & 8 & 117 & 0.892 & 7.447 & 1.9 \\
$\ZZ_p^*$ & 32 & 10 & 787 & 0.953 & 2.134 & 15.0 \\
$E(\FF_p)$ & 32 & 10 & 1\,558 & 1.310 & 2.156 & 16.7 \\
$\ZZ_p^*$ & 36 & 11 & 3\,039 & 1.277 & 0.901 & 56.4 \\
$E(\FF_p)$ & 36 & 11 & 6\,145 & 1.591 & 2.082 & 27.4 \\
\bottomrule
\end{tabular}
\end{table}

The build sits at $P\sim 0.8$--$1.6\,n^{2/3}$ and the table at
$n^{1/3}$ entries, close to the paper's $1.24\,n^{2/3}$ and
$n^{1/3}$. The online cost is $T\sim 1$--$7\,n^{1/3}$ against the
paper's $1.77\,n^{1/3}$. The $\sqrt{n}/T$ column is the point of the
method: already $7\times$--$56\times$ at these sizes, growing as
$n^{1/6}$, at the price of a precomputation larger than one
$\sqrt{n}$ search and worth it only amortised over many targets in a
fixed group. That is a statement about non-uniform security, not a
practical attack on a single logarithm.

A fixed-base doubling table, batched inversion, and a 16-byte
(fingerprint, exponent) table took the 36-bit curve build from $2.26$~s
to $0.44$~s on one core; four threads then build it in $0.11$~s. Class:
engineering.

\subsection{GLV folding}

On a CM curve the walk is folded by $\langle\psi\rangle$ (order $6$
for $j=0$, $4$ for $j=1728$) against the plain negation-map rho on the
same curve. The expected gain is $\sqrt{m/2}$ --- $1.73$ for $j=0$,
$1.41$ for $j=1728$. Measured over 20 instances: \texttt{glv-j0-32}
lands at $1.71\times$ ($S=0.988$ against $1.688$),
\texttt{glv-j1728-26} at $1.34\times$. The generic curve is identical
either way. Scatter is large because these are small curves; the
point is the folding, applied automatically once the $j$-invariant is
recognised.\art{docs/BENCHMARKS.md, ``GLV endomorphism-accelerated rho''}

\subsection{Cheon}

Cheon recovers $\alpha$ from $g,g^{\alpha},g^{\alpha^d}$ when $d$
divides $p-1$~\cite{cheon2006,cheon2010}. It is not bare ECDLP. With
$d\sim\sqrt{p}$ the library measures $1.9\times$ fewer group
operations than rho at 32~bits, $4.2\times$ at 40~bits, $18.6\times$
at 48~bits (each Cheon step is an exponentiation, $\sim 1.5\log_2 p$
operations).\art{docs/BENCHMARKS.md, ``Cheon's attack''} The companion
records a sharper statement on the $p-1$ and $p+1$ families with
fixed-base-comb BSGS: at 40~bits, $S=0.016$ on $p-1$ ($0.010\times$
matched rho, fitted exponent $0.231$ versus rho $0.487$); the $p+1$
family fits exponent $0.355$ versus rho $0.503$. These are advances
only against the DLP-with-auxiliary-inputs floor
$\Omega(\sqrt{p/d})$, not against the generic ECDLP
bound.\art{crypto:research/torsion_auxiliary_inputs_20260918/RESULTS.md}

\subsection{Cross-validation against the companion}
\label{sec:crossval}

Two independent implementations in the sibling recover the same
constants this library measures. They are cited as walk-checks. The
256-bit CUDA rho and the ECC2K-95 kernel have, per their READMEs, run
only through CPU emulation; no device-throughput number is taken from
them.

\paragraph{Rho $S=0.85$ versus BSGS $S=1.07$.}
On a 40-bit toy curve ($n=649\,523\,094\,257$), CPU emulation of the
companion kernels, 128 chains, every answer verified as $kG=Q$: rho
with the negation map measures $S=0.85$ (predicted
$0.96\times\sqrt{\pi n/4}$); rho without it $S=1.42$; BSGS with the
negation layout and batched $W=16$ measures $S=1.07$ (predicted
$1.00$), $1.26\times$ the rho reference and $1.07\times$ its own
floor. Textbook BSGS without negation is $S=1.44$ against a predicted
$1.41$. The batched stepper is engineering ($S$ unchanged, inversions
cut by $W$); the negation layout is the $\sqrt{2}$ the table
predicts; nothing here is an advance on rho's
floor.\art{crypto:gpu/ecc/README.md, ``Cost, measured''}

\paragraph{Frobenius-class walk, $0.94\times$ and $1.32\times$.}
The companion Koblitz kernel walks on classes of $\langle\tau,-1\rangle$
with the ECC2K-130 shape $P\leftarrow P+\tau^{j(P)}(P)$, $j=((g(x)/2)
\bmod 8)+3$, $g$ the normal-basis Hamming weight. On a 21-bit
subgroup, 128 walks, 24 trials: mean $251$ steps against a predicted
$\sqrt{\pi r/(4m)}=268$, i.e.\ $0.94\times$ the class-walk prediction
and $0.20\times$ a negation-only walk. On a 40-bit subgroup, 8 trials:
mean $191\,840$ against $145\,129$, i.e.\ $1.32\times$; $59$k of the
$192$k steps are the distinguished-point tail ($128$ walks $\times$
about $463$ steps to a point), and $145\mathrm{k}+59\mathrm{k}\approx
204\mathrm{k}$. Rho collision times have a standard deviation about
equal to their mean, so 24 and 8 trials resolve these to roughly
$\pm 20\%$ and $\pm 35\%$. They confirm the $\sqrt{2m}$ speedup; they
are too noisy to rank iteration
functions.\art{crypto:gpu/ecc2k/README.md, ``How this is verified''}

The same note records that dropping the $/2$ in $j$ collapses the walk
to four branches. Measured on the 40-bit curve, 64 walks, 160 trials:
$187\,831$ steps for the four-branch rule against $151\,005$ for the
eight-branch one. That is the Bailey et al.\ rule, and the one the
FPGA core of Section~\ref{sec:fpga-core} implements.
